%% paper_ko.tex — 학술지 형식 논문 (한국어 약 10,000단어)
%% 기업 지식 관리를 위한 적응형 검색 계획과 에이전트 기반 재검색
%%
%% 컴파일: xelatex paper_ko.tex && bibtex paper_ko && xelatex paper_ko.tex && xelatex paper_ko.tex

\documentclass[11pt, a4paper]{article}

%% ─── 한국어 지원 ───
\usepackage{kotex}
\usepackage{fontspec}
\setmainfont{AppleMyungjo}
\setsansfont{AppleGothic}
\setmonofont[Scale=0.88]{Courier New}

%% ─── 레이아웃 ───
\usepackage{microtype}
\usepackage[a4paper, top=25mm, bottom=25mm, left=30mm, right=28mm]{geometry}
\usepackage{amsmath}
\usepackage{amssymb}
\usepackage{booktabs}
\usepackage{multirow}
\usepackage{array}
\usepackage{tabularx}
\usepackage{graphicx}
\usepackage{float}
\usepackage{caption}
\captionsetup{font=small, labelfont=bf}
\usepackage[dvipsnames]{xcolor}
\usepackage[colorlinks=true, linkcolor=NavyBlue, citecolor=OliveGreen,
            urlcolor=MidnightBlue,
            pdftitle={기업 지식 관리를 위한 적응형 검색 계획과 에이전트 기반 재검색},
            pdfauthor={신원영}]{hyperref}
\usepackage[round,authoryear]{natbib}
\bibliographystyle{plainnat}
\usepackage{setspace}
\setstretch{1.2}
\usepackage{placeins}
\usepackage{enumitem}
\usepackage{titlesec}
\titleformat{\section}{\large\bfseries}{\thesection}{1em}{}
\titleformat{\subsection}{\normalsize\bfseries}{\thesubsection}{1em}{}
\titleformat{\subsubsection}{\normalsize\itshape}{\thesubsubsection}{1em}{}
\titlespacing*{\section}{0pt}{14pt}{6pt}
\titlespacing*{\subsection}{0pt}{10pt}{4pt}
\titlespacing*{\subsubsection}{0pt}{8pt}{4pt}

% ─────────────────────────────────────────────────────────────────────
\begin{document}

\begin{center}
{\Large\bfseries 기업 지식 관리를 위한 적응형 검색 계획과\\[0.2em]
에이전트 기반 재검색}\\[1.2em]
{\normalsize 신원영}\\
{\small 한양대학교 정보시스템학과}\\
{\small \texttt{pingu090@hanyang.ac.kr}}\\[0.8em]
{\small 2026년 6월}
\end{center}

% ─────────────────────────────────────────────────────────────────────
\begin{abstract}
기업 지식 검색은 일반 정보 검색 연구에서 다루지 않는 고유한 과제를 제시한다. 단일 문서 코퍼스에 API 기술 참조서, 인사 정책 매뉴얼, 보안 컴플라이언스 가이드, 프로젝트 회의록이 혼재하며, 각각 효과적인 접근을 위해 서로 다른 검색 신호를 필요로 한다. 본 논문은 이 문제에 대한 두 가지 상호보완적 기여를 제시한다. 첫째, 합성 기업 지식 데이터셋(SEKD: Synthetic Enterprise Knowledge Dataset)을 구축한다. SEKD는 6가지 추론 유형과 8가지 난이도 요인에 걸쳐 405개의 쿼리를 포함하는 완전 주석 처리된 120개 문서, 1,206개 청크 코퍼스로, 청크 수준 정답 주석과 오라클 검색 전략 레이블을 함께 제공한다. SEKD를 사용하여 기업 지식 쿼리에 대한 5가지 검색 시스템의 최초 통제된 체계적 비교를 수행한다: BM25 희소 검색, BGE-small-en-v1.5를 활용한 밀집 검색, 하이브리드 역순위 융합(RRF), 7개 규칙 기반 적응형 플래너, 제로샷 GPT-5 전략 계획. 실험 결과, 하이브리드 RRF가 모든 주요 지표에서 단일 전략 기준선을 압도하고(Recall@10=0.745, MRR=0.530, nDCG@10=0.566), 7개 규칙 결정론적 플래너가 운영 비용 없이 전체 시스템 중 최고 Recall@10(0.755)을 달성하며, GPT-5가 380개 쿼리당 \$6.18의 비용으로 순위 품질 향상(MRR +2.0\%, nDCG@10 +1.2\%)을 제공함을 확인하였다. 결정적으로, 모든 평가된 시스템이 예외 절 쿼리(Recall@10~$\leq$~0.412)와 시계열 추론 쿼리(Recall@10~$\leq$~0.200)에서 체계적으로 실패한다. 둘째, 이러한 지속적 실패에 동기를 받아 V2 에이전트 기반 RAG 파이프라인을 도입한다. 이는 반복적 검증--재검색 루프로 V1 검색 인프라를 감싸는 5-에이전트 LangGraph 시스템이다. 검증 에이전트는 검색된 청크를 유형별 관련성 임계값과 비교하여 점수를 매기고, 거부 시 실패 원인을 진단하며, 수정적 재검색을 위한 전략을 제안한다. 4개의 통제된 제거 실험(E1--E4)을 설계하여 재검색 루프, 검증 점수화 품질, 쿼리 분류 정확도의 기여를 분리한다. 결과는 어떤 에이전트 구성 요소가 정적 계획을 넘어 측정 가능한 가치를 더하는지, 그리고 어떤 실패 모드가 반복적 수정 후에도 여전히 불가능한지를 보여준다.
\end{abstract}

\noindent\textbf{키워드:} 기업 검색, 하이브리드 RRF, 적응형 검색 계획, 에이전트 기반 RAG, 재검색 루프, LangGraph, BM25, 밀집 구문 검색, 기업 지식 관리.

% ─────────────────────────────────────────────────────────────────────
\section{서론}
\label{sec:intro}

현대 기업은 대량의 이질적인 구조화 및 반구조화 지식을 생성하고 유지한다. 전형적인 중간 규모 기술 조직은 API 기술 참조서, 인사 정책 매뉴얼, 시스템 아키텍처 문서, 프로젝트 회의록, 보안 컴플라이언스 가이드, 출장비 정책 등의 운영 정책 문서에 대한 별도 문서 저장소를 유지한다. 이러한 코퍼스가 규모와 다양성 면에서 성장함에 따라, 관련 정보를 효율적으로 찾는 것이 일차적인 조직 문제가 된다. 수백 개의 정책 문서를 검색하는 직원, 시스템 설계 아카이브를 조회하는 엔지니어, 자동화된 컴플라이언스 확인 시스템 모두 올바른 정보를 안정적으로 제공하는 검색 시스템에 의존한다.

검색 증강 생성(RAG: Retrieval-Augmented Generation)은 지식 집약적 자연어 처리의 지배적 패러다임으로 부상하였다 \citep{lewis2020rag}. RAG 시스템에서 업스트림 검색 구성 요소는 코퍼스에서 후보 문서나 구문을 검색하고, 다운스트림 언어 모델은 검색된 구문을 조건으로 출력을 생성한다. 따라서 검색 단계가 핵심 병목이다: 관련 문서를 검색하지 못하면, 어떤 수준의 다운스트림 생성 품질도 이를 보완할 수 없다. 검색 품질은 두 가지 주요 차원으로 측정된다---\emph{커버리지}(시스템이 모든 관련 문서를 찾는가?)와 \emph{순위 품질}(가장 관련성 높은 문서를 최상위에 배치하는가?).

기업 코퍼스는 IR 연구를 지배하는 웹 검색 및 공개 도메인 질문 답변 설정과 구별하는 세 가지 구조적 과제를 제시한다. 첫째, \emph{쿼리 이질성}: 기업 쿼리는 단순한 단일 도약 사실 조회(``호텔 환급 한도는 얼마인가?'')부터 복잡한 다중 도약 추론(``데이터 내보내기 서비스와 동일한 인증 방법을 사용하는 API 엔드포인트는?''), 시계열 순서 쿼리(``2024년 3분기에 유효한 정책 버전은?''), 정책 예외 절 추출(``직원이 사전 승인을 면제받는 조건은?''), 다중 속성 집계 작업(``이중 인증을 의무화하는 모든 보안 섹션 나열'') 등 광범위한 추론 유형에 걸쳐 있다. 서로 다른 추론 유형은 근본적으로 다른 검색 신호로부터 이익을 얻는다.

둘째, \emph{코퍼스 이질성}: 동일한 검색 방법이 API 참조서(식별자가 풍부하고 구조가 무거움)와 HR 정책(의미론이 무겁고 예외가 밀집함)에 걸쳐 최적일 수 없다. 의미적 유사성 매칭을 위해 미세 조정된 밀집 인코더는 패러프레이징된 정책 쿼리에서 탁월할 수 있지만, 정확한 버전 번호나 엔드포인트 경로를 매칭해야 하는 쿼리에서 실패할 수 있다.

셋째, \emph{단일 패스 경직성}: 모든 표준 검색 아키텍처는 쿼리당 전략 선택 결정을 한 번만 내리며, 초기 결과가 불충분할 경우 복구할 메커니즘이 없다. 이 경직성은 쿼리 의미와 검색 신호 사이의 불일치가 구조적인 예외 및 시계열 쿼리 유형에 특히 비용이 많이 든다.

\emph{적응형 검색 계획}은 각 쿼리를 해당 특성에 가장 적합한 검색 백엔드로 동적으로 라우팅함으로써 두 번째 과제를 해결한다. 그러나 적응형 계획만으로는 세 번째 과제를 해결할 수 없다: 플래너의 초기 전략 선택이 잘못되거나, 전략에 관계없이 초기 검색이 불충분하면 수정 메커니즘이 존재하지 않는다.

본 논문은 기업 지식 검색의 솔루션으로 적응형 계획과 에이전트 기반 재검색 모두를 조사한다. 연구를 13개의 연구 질문으로 구성한다. RQ1--RQ3은 검색 기준선 환경을 수립한다: 하이브리드 RRF가 단일 방법 기준선을 능가하는가; 문서 범주별 최적 방법은 무엇인가; 추론 유형별 성능은 어떻게 변하는가? RQ4--RQ7은 적응형 계획을 평가한다: 규칙 기반 계획이 고정 전략 검색을 능가할 수 있는가; 최적의 LLM 프롬프트 전략은 무엇인가; GPT-5가 규칙 플래너를 능가하는가; LLM 계획의 실패 모드는 무엇인가? RQ8--RQ13은 V2 에이전트 확장을 평가한다: 파이프라인이 V1 플래너를 능가하는가; 각 에이전트 구성 요소의 기여는 무엇인가; 재검색이 예외 및 시계열 실패를 해결하는가; 어떤 실패 모드가 남는가?

실험 프로그램은 두 단계로 진행된다. \emph{V1 평가}(3--5절)는 SEKD에서 5가지 검색 시스템 모두의 성능을 특성화하고, 단일 패스 검색의 이상적인 조건 하에서 성능 상한을 수립하며, 어떤 정적 계획 메커니즘도 해결할 수 없는 실패 모드를 식별한다. \emph{V2 평가}(4절)는 에이전트 파이프라인 아키텍처와 4개의 통제된 제거 실험(E1--E4)을 도입하며, 검증 에이전트의 임계값 보정 문제를 드러내는 경험적 결과를 제시한다.

본 논문의 기여는 다음과 같다. \textbf{(1) 데이터셋}: 검색 방법, 플래너, 에이전트 시스템의 체계적 평가를 지원하는 405개 주석 쿼리를 포함하는 통제된 기업 코퍼스 SEKD를 구축하고 공개한다. \textbf{(2) V1 기준선 연구}: 기업 지식 쿼리에 대한 5가지 검색 시스템의 최초 체계적 비교를 보고하여, 추론 유형 및 문서 범주에 걸쳐 경험적으로 근거한 성능 벤치마크를 수립한다. \textbf{(3) V2 에이전트 시스템}: V1 계획 구성 요소를 폴백으로 보존하면서 V1에서 식별된 예외 및 시계열 실패 모드를 목표로 하는 반복적 검증--재검색을 포함하는 5-에이전트 LangGraph 파이프라인을 구현한다. \textbf{(4) 제거 연구}: 통제된 실험을 통해 재검색, 적응형 검증 점수화, 쿼리 분류 정확도의 기여를 분리한다. \textbf{(5) 실용적 권고사항}: 전체 실험 프로그램을 기반으로 기업 RAG 실무자를 위한 정량적이고 실행 가능한 설계 지침을 제공한다.

% ─────────────────────────────────────────────────────────────────────
\section{관련 연구}
\label{sec:related}

\subsection{희소 검색과 밀집 검색}

BM25 \citep{robertson2009bm25,robertson1995okapi} 계열의 희소 검색 모델은 포화 및 정규화 매개변수를 포함하는 단어 빈도와 역문서 빈도로 문서에 점수를 매긴다. Okapi BM25 변형(수식 \ref{eq:bm25})은 고빈도 단어가 점수를 지배하는 것을 방지하는 매개변수 $k_1$을 포함하는 준선형 TF 포화 함수를 적용하고, 더 긴 문서의 이점을 조정하는 문서 길이 정규화 항을 포함한다.

\begin{equation}
\text{BM25}(q,d) = \sum_{t \in q} \text{IDF}(t) \cdot \frac{f_{t,d} \cdot (k_1 + 1)}{f_{t,d} + k_1 \cdot (1 - b + b \cdot |d| / \text{avgdl})}
\label{eq:bm25}
\end{equation}

여기서 $f_{t,d}$는 문서 $d$에서 단어 $t$의 빈도, $|d|$는 문서 길이, $\text{avgdl}$은 코퍼스 평균 문서 길이, $\text{IDF}(t) = \log \frac{N - n_t + 0.5}{n_t + 0.5}$이며 $N$은 코퍼스 크기, $n_t$는 문서 빈도이다. 본 구현에서는 $k_1{=}1.5$, $b{=}0.75$를 표준 기업 검색 구성을 따라 적용한다. 밀집 구문 검색(DPR) \citep{karpukhin2020dpr}은 질문 및 구문 인코더를 별도로 학습하고 근사 최근접 이웃 검색을 통해 검색하는 양방향 인코더 밀집 검색을 도입하였다. \citet{xiong2020approximate}는 FAISS 기반 근사 검색을 사용하여 밀집 검색을 확장하였다. 본 연구에서 사용하는 BAAI/bge-small-en-v1.5는 낮은 계산 비용으로 경쟁력 있는 BEIR 벤치마크 성능을 달성하는 소형(384차원) 양방향 인코더이다 \citep{thakur2021beir}.

\subsection{하이브리드 검색과 융합}

\citet{cormack2009rrf}는 역순위 점수를 누적하여 여러 순위 목록을 결합하는 매개변수 없는 융합 방법인 역순위 융합(RRF)을 도입하였다. RRF의 주요 장점은 학습이 필요 없고, 검색 시스템 간 점수 크기 차이에 강건하며, 한 순위 목록에 문서가 없어도 허용한다는 것이다. \citet{luan2021sparse_dense}는 희소-밀집 하이브리드 검색이 어느 단일 모드 접근법보다 일관되게 우수하며, 특히 상호보완적인 어휘 및 의미적 요구사항을 가진 쿼리에서 그러함을 입증하였다. 본 논문에서는 SEKD 기업 쿼리 분포에서 상보성과 RRF 융합 이점 모두를 확인한다.

\subsection{적응형 검색 계획 및 쿼리 라우팅}

쿼리 특성에 따른 검색 방법 또는 모델 선택(쿼리 라우팅)은 쿼리 분류 \citep{salton1975smart}, 선택적 검색, 쿼리 적응형 융합 등의 레이블 하에 IR 문헌에서 연구되었다. \citet{nogueira2019passage}는 쿼리 유형이 검색 성능을 조절하며 쿼리 특성이 어떤 검색 방법이 가장 효과적일지 예측할 수 있음을 보였다. LLM 기반 쿼리 분석은 \citet{ma2023query}가 탐구하였으며, 명시적 구조화된 메타데이터가 프롬프트의 모델 크기보다 전략 분류 정확도에 더 중요함을 발견한 본 GPT-5 평가로 이 연구를 확장한다. \citet{shi2023large}는 LLM 프롬프트에 주입된 관련 없는 맥락이 추론 품질을 실질적으로 저하시킬 수 있음을 입증하여, 계획 작업에서 신중한 프롬프트 엔지니어링의 필요성을 강조한다.

\subsection{기업 정보 검색}

기업 정보 검색은 웹 검색이나 공개 도메인 질문 답변에 비해 학술 문헌에서 비교적 적은 주목을 받아왔지만, 실용적 중요성은 크다. 기업 환경은 여러 구조적 측면에서 웹 IR과 다르다: 문서 코퍼스는 유한하고 큐레이션되어 있으며, 문서는 설계상 이질적이고, 쿼리는 종종 고위험이다(컴플라이언스 검증, 법적 연구, 기술 사양 조회). TREC 기업 트랙 \citep{voorhees2005trec}은 기업 코퍼스에서 전문가 검색 및 알려진 항목 검색을 위한 벤치마크를 수립하였다. BEIR \citep{thakur2021beir}은 기업 인접 도메인을 포함한 크로스 도메인 제로샷 검색 벤치마크를 제공하였다. 본 연구는 범주 혼합, 추론 유형 분포, 난이도 요인 프로필을 포함하여 기업 지식 코퍼스의 구조적 특성을 캡처하도록 특별히 설계된 코퍼스를 기여한다.

\subsection{에이전트 기반 및 반복적 RAG}

ReAct 프레임워크 \citep{yao2023react}는 내부 연쇄 추론 추론이 안내하는 반복적 도구 사용을 가능하게 하는 LLM 에이전트를 위한 추론-행동 교차 패러다임을 도입하였다. FLARE \citep{ji2023flare}는 LLM의 생성 신뢰도(토큰 확률로 측정)가 임계값 아래로 떨어질 때 추가 맥락 검색을 동적으로 결정한다. Self-RAG \citep{asai2024selfrag}는 검색이 필요한지, 검색된 구문이 관련성 있는지, 생성된 응답이 검색된 증거에 의해 지지되는지를 나타내는 특별 비평 토큰을 예측하도록 모델을 학습시킨다. MuSiQue \citep{trivedi2022musique}는 특정 순서로 여러 문서에서 증거를 필요로 하는 구성적 다중 도약 질문 답변을 도입하였다.

본 V2 시스템은 Self-RAG 및 FLARE와 핵심 설계 차원에서 다르다: LLM 생성 내에 검색 결정을 내장하는 대신, 검색 검증을 생성 전에 작동하는 독립적인 점수화 함수로 구현하여, 검색 품질 평가와 답변 합성 사이의 깔끔한 분리를 가능하게 한다. LangGraph \citep{langgraph2024}는 조건부 엣지(재검색 역방향 엣지), 에이전트 노드 간 공유 타입 상태, 개별 구성 요소를 독립적으로 교체하거나 제거할 수 있는 에이전트 수준 모듈성을 지원하는 V2 에이전트 구현을 위한 상태형 방향 그래프 프레임워크를 제공한다.

% ─────────────────────────────────────────────────────────────────────
\section{데이터셋: 합성 기업 지식 데이터셋(SEKD)}
\label{sec:dataset}

\subsection{코퍼스 구성}

SEKD 코퍼스는 6가지 기업 문서 범주에 걸쳐 120개 문서로 구성된다: (1)~\textit{API 문서}---엔드포인트 사양, 인증 방식, 속도 제한, 통합 가이드를 설명하는 기술 참조 문서; (2)~\textit{인사 정책}---휴가 자격, 성과 검토 프로세스, 보상 지침, 징계 절차를 다루는 고용 정책 문서; (3)~\textit{프로젝트 회의록}---결정 사항, 액션 항목, 프로젝트 현황 업데이트, 종속성 추적이 포함된 구조화된 회의 의사록; (4)~\textit{보안 정책}---인증 요구사항, 데이터 처리 의무, 접근 제어 사양, 인시던트 대응 절차를 다루는 컴플라이언스 및 보안 거버넌스 문서; (5)~\textit{시스템 설계 문서}---시스템 구성 요소, 데이터 흐름, 통합 패턴, 배포 구성을 설명하는 아키텍처 사양; (6)~\textit{출장 정책}---환급 한도, 승인 워크플로우, 예약 절차, 비용 범주를 명시하는 운영 정책 문서.

각 문서는 무작위화된 엔터티 값(프로젝트 이름, 직원 ID, API 엔드포인트 경로, 정책 섹션 코드, 금전적 임계값, 날짜 범위)을 포함하는 Jinja2 템플릿 파이프라인을 사용하여 생성되었다. 무작위화는 재현성을 위해 고정된 값으로 시드되었으며, 사소한 키워드 매칭을 방지하기에 충분한 다양성을 생성하였다. 다중 도약 쿼리 유형을 위해 문서 간 일관성이 의도적으로 포함되어, 시스템 설계 문서와 해당 API 참조서가 다중 도약 쿼리가 두 문서 모두에서 정보를 합성해야 할 수 있도록 엔터티 이름을 공유한다.

결과 코퍼스에는 1,206개의 섹션 인식 청크가 포함된다. 청킹은 3단계 전략을 따른다: (1)~보존된 섹션 메타데이터(문서 ID, 범주, 섹션 경로, 섹션 번호)와 함께 섹션 헤더에서 분할; (2)~50토큰 오버랩으로 단락 경계에서 과대 섹션을 재귀적으로 분할하여 문맥 경계 정보 손실 방지; (3)~너무 짧은 터미널 섹션(50토큰 미만)을 이전 형제 섹션과 병합하여 사소하게 짧은 청크가 검색 결과에 나타나지 않도록 방지. 평균 청크 길이는 248토큰이며, 사분위수 범위는 156--327토큰으로 문서 범주에 걸쳐 섹션 길이의 다양성을 반영한다.

\subsection{쿼리 세트 구성 및 주석}

405개 쿼리 세트는 3단계 프로세스를 통해 구성되었다. \emph{생성 단계}에서 쿼리 템플릿이 각 문서에 대해 인스턴스화되었으며, 쿼리는 모든 추론 유형과 난이도 요인을 포함하도록 설계되었다. \emph{주석 단계}에서 직접 인용 증거 추출 절차가 각 답변의 완전하고 분해할 수 없는 사실적 근거를 포함하는 SEKD 청크의 최소 세트를 식별하여 쿼리당 정답 청크 ID 세트를 생성하였다. \emph{품질 단계}에서 모호한 증거를 가진 쿼리, SEKD에서 이용할 수 없는 세계 지식이 필요한 쿼리, 구조적 중복이 제거되었다. 최종 세트에는 380개의 답변 가능 쿼리(모든 평가에 사용)와 25개의 답변 불가능 쿼리가 포함된다.

추론 유형 분포: 단일\_도약(220, 57.9\%), 다중\_도약(53, 13.9\%), 집계(35, 9.2\%), 비교(33, 8.7\%), 예외(34, 8.9\%), 시계열(5, 1.3\%). 단일\_도약 지배는 기업 쿼리의 현실적인 분포를 반영한다: 대부분의 조회 쿼리는 하나의 문서 청크를 필요로 하는 직접적인 단일 도약 검색이다. 오라클 검색 전략 레이블은 검색 신호 분석에 근거한 휴리스틱 매핑으로 할당되었다: 시계열→bm25, 예외→dense, 다중\_도약→hybrid, 집계→hybrid, 비교→hybrid, 단일\_도약→쿼리 내용 및 범주에 따른 전략. 오라클 레이블 분포: Hybrid(249, 65.5\%), Dense(69, 18.2\%), BM25(62, 16.3\%).

\subsection{난이도 요인 및 주석 프로토콜}

8가지 난이도 요인이 쿼리 복잡성을 주석 처리한다: \emph{시계열\_종속성}, \emph{예외\_절}, \emph{다중\_문서\_종속성}, \emph{표\_종속성}, \emph{문서\_버전\_충돌}, \emph{부정}, \emph{대명사\_해석}, \emph{암묵적\_추론}. 각 주석자는 직접 인용 증거 추출로 관련 청크 세트를 식별하였다. 쿼리는 코퍼스에 답변에 충분한 증거를 포함하는 청크가 없거나, 답변이 SEKD에 없는 세계 지식을 필요로 하거나, 증거 범위가 여러 충돌하는 청크에 걸쳐 모호한 경우 답변 불가능으로 표시되었다. 답변 가능 세트(380개 쿼리)의 난이도 요인 분포: 시계열\_종속성(5, 1.3\%), 예외\_절(34, 8.9\%), 다중\_문서\_종속성(53, 13.9\%), 표\_종속성(41, 10.8\%), 문서\_버전\_충돌(18, 4.7\%), 부정(29, 7.6\%), 대명사\_해석(22, 5.8\%), 암묵적\_추론(67, 17.6\%).

\subsection{평가 프로토콜}

검색 평가는 각 답변 가능 쿼리를 청크 ID에 대한 재현율-정밀도 작업으로 처리한다. 청크는 검색 시스템이 반환한 상위 $k$ 결과에 나타나면 \emph{검색됨}으로, 해당 쿼리의 정답 청크 ID 세트의 구성원이면 \emph{관련됨}으로 처리된다. 표준 지표인 Recall@K, MRR, nDCG@K, Hit@K가 $k \in \{1, 3, 5, 10\}$에서 계산된다. 모든 시스템은 직접적인 지표 비교를 가능하게 하는 동일한 380개 답변 가능 쿼리와 동일한 코퍼스에서 평가된다.

기본 평가 결과는 검색된 구문에 LLM 생성을 조건화하는 RAG 파이프라인을 위한 현실적인 검색 예산을 반영하는 $k=10$에서 보고된다. MRR과 nDCG@10은 커버리지와 순위 품질을 결합하며, nDCG@10은 상위 10개 창 내에서 관련 결과의 위치에 더 많은 가중치를 부여한다.

% ─────────────────────────────────────────────────────────────────────
\section{V1 검색 시스템}
\label{sec:v1}

\subsection{검색 백엔드}

\textbf{BM25.} 표준 Okapi 점수화 함수($k_1{=}1.5$, $b{=}0.75$, $\epsilon{=}0.25$)를 사용하여 BM25를 구현한다. BM25는 정확한 토큰 매칭으로부터 이익을 얻으며, 정확한 식별자(정책 코드, 버전 번호, API 엔드포인트 경로, 금전적 임계값)를 포함하는 쿼리에 효과적이다.

\textbf{밀집 검색.} BAAI/bge-small-en-v1.5(384차원)를 사용하여 문서와 쿼리를 인코딩한다. 문서는 오프라인에서 인코딩되어 8개 클러스터가 있는 FAISS IVF 인덱스에 저장된다. 밀집 검색은 의미적 매칭으로부터 이익을 얻지만, 날짜 토큰이 정확한 순차적 의미를 전달하는 시계열 쿼리에서 임베딩 표현이 비슷하게 생성되어 제로 검색 재현율을 초래한다.

\textbf{하이브리드 RRF.} $k{=}10$으로 역순위 융합 \citep{cormack2009rrf}을 사용하여 BM25와 밀집 순위 목록을 융합한다. 청크 $c$에 대한 융합 점수는 $\text{RRF}(c) = \sum_{i \in \{\text{bm25,dense}\}} (k + r_i(c))^{-1}$이며, $r_i(c)$는 목록 $i$에서 청크 $c$의 순위이다(없으면 $\infty$). RRF는 매개변수가 없고, 누락된 청크를 우아하게 처리하며, 검색 시스템 간 점수 크기 차이에 강건하다.

\subsection{규칙 기반 적응형 플래너}

규칙 플래너는 우선순위 순서로 7개의 결정론적 규칙을 평가한다(표 \ref{tab:rules}). 규칙 R01--R04는 최적 전략이 명확한 추론 유형을 대상으로 한다. R07은 무조건적인 Hybrid 기본값이다.

\begin{table}[htbp]
\centering
\caption{규칙 플래너 결정 논리: 우선순위 순서로 7개 규칙}
\label{tab:rules}
\begin{tabular}{lllc}
\toprule
\textbf{규칙} & \textbf{조건} & \textbf{전략} & \textbf{SSA} \\
\midrule
R01 & reasoning\_type $=$ temporal          & BM25   & 100\% \\
R02 & reasoning\_type $=$ exception         & Dense  & 100\% \\
R03 & reasoning\_type $=$ multi\_hop        & Hybrid & 100\% \\
R04 & reasoning\_type $=$ aggregation       & Hybrid & 100\% \\
R05 & difficulty $=$ document\_version\_conflict & BM25 & N/A \\
R06 & difficulty $=$ table\_dependency      & Dense  & N/A  \\
R07 & (매칭되지 않는 모든 쿼리에 대한 기본값) & Hybrid & 83.9\% \\
\midrule
\multicolumn{3}{l}{\textbf{전체 SSA}} & \textbf{78.4\%} \\
\bottomrule
\end{tabular}
\end{table}

플래너의 주요 약점은 BM25 과소 예측이다: API 문서 및 출장 정책 범주의 62개 오라클-BM25 단일 도약 쿼리 중 5개만 BM25로 라우팅한다. R07의 Hybrid 기본값이 단일 도약 쿼리에 대해 지배적이다.

\subsection{범주별 검색 성능}

표 \ref{tab:category}는 3개의 검색 백엔드와 규칙 플래너에 대한 문서 범주별 Recall@10을 보고한다. 범주 수준 분석은 문서 구조가 다른 어떤 단일 요인보다 검색 난이도를 더 많이 결정함을 보여준다.

\begin{table}[htbp]
\centering
\caption{문서 범주별 Recall@10. 범주별 최고 성능은 볼드체.}
\label{tab:category}
\begin{tabular}{lrccccc}
\toprule
\textbf{범주} & \textbf{n} & \textbf{BM25} & \textbf{밀집} & \textbf{하이브리드} & \textbf{규칙} & \textbf{GPT-5} \\
\midrule
API 문서      & 68  & 0.734 & 0.751 & 0.755 & \textbf{0.773} & 0.763 \\
인사 정책     & 54  & 0.649 & 0.701 & 0.702 & 0.712 & \textbf{0.718} \\
회의록        & 100 & 0.772 & 0.784 & \textbf{0.793} & 0.787 & 0.783 \\
보안 정책     & 64  & 0.696 & 0.721 & 0.729 & 0.741 & \textbf{0.747} \\
시스템 설계   & 45  & 0.703 & \textbf{0.756} & 0.752 & 0.748 & 0.740 \\
출장 정책     & 49  & 0.693 & 0.718 & 0.727 & \textbf{0.744} & 0.741 \\
\midrule
\textbf{전체} & 380 & 0.713 & 0.740 & 0.745 & \textbf{0.755} & 0.750 \\
\bottomrule
\end{tabular}
\end{table}

\textbf{API 문서} 쿼리는 시스템 간 가장 큰 절대 분산을 보인다(BM25에서 규칙 플래너까지 +3.9 pp). API 문서는 정확한 식별자(엔드포인트 경로, 버전 문자열, HTTP 메서드 이름)가 풍부하여 BM25를 강력한 경쟁자로 만들고 규칙 플래너의 대상 BM25 라우팅(R01, R05)을 특히 효과적으로 만든다. \textbf{시스템 설계}는 Dense가 Hybrid를 능가하는 유일한 범주이며(+0.4 pp), 이 범주에서는 정확한 식별자 조회보다 의미적 매칭이 지배적임을 시사한다.

\subsection{GPT-5 LLM 플래너}

구조화된 JSON 출력 강제를 포함한 OpenAI Chat Completions API를 통해 제로샷 검색 전략 플래너로 GPT-5를 평가한다. 프롬프트는 명시적 전략 지침(``시계열~$\to$~bm25'', ``예외~$\to$~dense'' 등)을 인코딩하고 구조화된 형식으로 쿼리 메타데이터를 제공한다. 세 가지 프롬프트 변형에 대한 예비 연구에서: 제로샷 자연어(SSA 69.5\%), 메타데이터가 있는 구조화(SSA 82.1\%), 5개 전략 예제가 있는 퓨샷(SSA 81.8\%). 구조화 프롬프트가 퓨샷 토큰 비용의 약 절반으로 최고 SSA를 달성하여 주요 평가에 선택되었다. 모든 380개 쿼리가 유효한 JSON을 반환했다(파싱 실패율: 0.0\%). 총 비용: \$6.18(\$0.016/쿼리); 평균 지연: 5,649~ms.

% ─────────────────────────────────────────────────────────────────────
\section{V2 에이전트 기반 RAG}
\label{sec:v2}

\subsection{설계 원칙 및 동기}

V1 평가는 두 가지 제한 사항 클래스를 드러낸다. \emph{클래스~1(구조적 신호 불일치):} 예외 및 시계열 쿼리는 이용 가능한 검색 신호가 필요한 쿼리 의미를 캡처하지 못하기 때문에 실패한다. 이 실패는 전략 선택만으로는 수정될 수 없다. \emph{클래스~2(단일 패스 경직성):} 초기 검색이 불충분할 때, 이유를 진단하거나 수정 전략을 시도할 메커니즘이 없다.

V2는 세 가지 안내 설계 원칙을 통해 클래스~2 실패를 직접 해결한다: \textbf{(P1) 단순 라우팅이 아닌 추론 분해}: 복잡한 쿼리는 다중 도약 쿼리를 위한 순차적 증거 수집을 가능하게 하는 하위 쿼리별 검색 계획과 함께 순서화된 하위 쿼리로 분해된다. \textbf{(P2) 단순 검색이 아닌 증거 검증}: 검색된 청크는 신호 일치 관련성 점수화기를 사용하여 특정 하위 쿼리 텍스트와 비교하여 점수가 매겨지고, 점수화된 세트가 증거 요구사항을 충족하지 못하면 수정 전략과 함께 재검색이 트리거된다. \textbf{(P3) V1을 폴백으로 보존}: V1 구성 요소(규칙 플래너, BM25/Dense/Hybrid 인덱스)가 V2의 핵심으로 유지된다.

\subsection{에이전트 아키텍처}

V2 파이프라인은 타입화된 \texttt{AgentState} 객체를 공유하는 5개 에이전트 노드로 구성된 LangGraph StateGraph \citep{langgraph2024}로 구현된다. 기본 실행 흐름은:

\begin{equation*}
\underbrace{\text{QAA}}_{\text{쿼리 분석}} \to
\underbrace{\text{PA}}_{\text{계획}} \to
\underbrace{\text{RA}}_{\text{검색}} \to
\underbrace{\text{VA}}_{\text{검증}}
\xrightarrow{\text{통과}} \underbrace{\text{AA}}_{\text{답변}} \to \text{종료}
\end{equation*}

\noindent 검증 실패 시 재검색 역방향 엣지:

\begin{equation*}
\text{VA} \xrightarrow{\text{실패},\ \texttt{loop\_count} < 3} \text{루프카운터} \to \text{PA}
\end{equation*}

상태 설계는 \texttt{retrieval\_outputs} 필드가 루프 전반의 결과를 누적하지 않고 가장 최근의 검색 시도만 보유하는 방식으로, 검색 지표가 모든 시도의 합집합이 아닌 파이프라인의 최상의 전달 결과를 반영하도록 의도적으로 상태 비저장 방식을 따른다.

\subsubsection*{쿼리 분석 에이전트(QAA).}

QAA는 수신 쿼리를 6가지 추론 유형 중 하나로 분류하고 난이도 신호를 추출한다. 세 가지 운영 모드가 지원된다: \textit{오라클 모드}(실험 E1--E3): 지면 진실 추론 유형이 쿼리 메타데이터에서 직접 제공된다. \textit{휴리스틱 모드}(실험 E4): 6패턴 정규식 분류기가 표면 특성에 기반하여 추론 유형을 할당한다. \textit{LLM 모드}(운영 배포): 구조화된 JSON 프롬프트가 LLM 제공자에게 추론 유형, 난이도 신호, 분해 플래그, 하위 쿼리를 요청한다.

\subsubsection*{계획 에이전트(PA).}

PA는 3단계 결정 프로세스를 통해 각 하위 쿼리를 검색 전략에 매핑한다. 첫째, \texttt{loop\_count}~$>0$이고 VA가 \texttt{suggested\_strategy}를 제공했다면, 해당 전략이 무조건적 우선순위를 가진다. 둘째, 첫 번째 검색 시도(루프 카운트 0)에서는 V1 규칙 플래너가 적용된다. 셋째, 매칭 규칙의 신뢰도가 0.75 미만이고 LLM 제공자가 이용 가능하면 LLM 계획 호출이 규칙 결과를 보완한다.

\subsubsection*{검색 에이전트(RA).}

RA는 사전 구축된 SEKD 인덱스에 대한 세 가지 MCP 도구 래퍼(\texttt{bm25\_search}, \texttt{dense\_search}, \texttt{hybrid\_search})를 사용하여 검색 계획을 실행한다. 실행 모드는 QueryPlan의 추론 유형으로 결정된다: 다중\_도약의 경우 \textit{순차적}(이전 하위 쿼리 결과에서 슬롯 채우기), 집계의 경우 \textit{병렬-병합}, 다른 모든 유형의 경우 \textit{독립적}.

\subsubsection*{검증 에이전트(VA).}

VA는 V2 파이프라인의 핵심 피드백 메커니즘이다. 수식 (\ref{eq:pass})을 검색된 청크 세트에 적용하고, 통과 또는 실패 결정을 생성하며, 실패 시 수정 전략을 포함하는 구조화된 실패 진단을 출력한다.

\begin{equation}
\text{통과} \iff |\mathcal{R}| \geq \delta, \quad
|\mathcal{R}| = \bigl|\{c \in \mathcal{C} : \text{score}(c,\, q_c) \geq \theta_{\tau_c}\}\bigr|
\label{eq:pass}
\end{equation}

여기서 $\mathcal{C}$는 검색된 청크 세트, $q_c$는 청크 $c$에 연관된 하위 쿼리 텍스트, $\tau_c$는 추론 유형, $\theta_{\tau}$는 유형별 임계값, $\delta=3$은 최소 통과 청크 수이다. 하위 유형 임계값은 경험적 점수 분포를 반영한다: 기본 $\theta=0.30$; 예외, 다중\_도약, 시계열: $\theta=0.20$--$0.22$.

5가지 교환 가능한 관련성 점수화기가 구현된다(표 \ref{tab:scorers}). \texttt{adaptive} 점수화기(E1, E2, E4의 기본값)는 \texttt{strategy\_used} 필드에 기반하여 각 청크를 BM25 또는 Dense 점수화기로 라우팅한다. 순위 프록시 점수화기(E3)는 쿼리별 관련성이 아닌 검색 순위에만 의존한다.

\begin{table}[htbp]
\centering
\caption{검증 에이전트 관련성 점수화기}
\label{tab:scorers}
\begin{tabular}{lll}
\toprule
\textbf{점수화기} & \textbf{신호} & \textbf{비용} \\
\midrule
\texttt{rank\_proxy}  & 순위 $r$에 대해 $\max(0,\ 1 - (r{-}1) \cdot 0.1)$ & 없음 \\
\texttt{bm25}         & 정규화된 BM25 점수 (청크 대 하위 쿼리) & 무시 가능 \\
\texttt{dense}        & 인덱스의 코사인 유사도 (\texttt{chunk.score}) & 없음 \\
\texttt{adaptive}     & \texttt{strategy\_used}에 따라 bm25/dense/평균으로 라우팅 & 무시 가능 \\
\texttt{llm}          & GPT-4o-mini 배치 점수화 (5청크/호출) & $\sim$\$0.001/10청크 \\
\bottomrule
\end{tabular}
\end{table}

VA의 실패 진단 함수는 검증이 실패하면 점수 분포를 분석하여 구조화된 진단(표 \ref{tab:failure_diag})과 다음 루프 반복에서 PA가 사용할 \texttt{suggested\_strategy}를 생성한다.

\begin{table}[htbp]
\centering
\caption{검증 에이전트 실패 진단 논리}
\label{tab:failure_diag}
\begin{tabular}{lll}
\toprule
\textbf{진단} & \textbf{조건} & \textbf{권고사항} \\
\midrule
\texttt{no\_chunks}             & 검색된 청크 없음               & Hybrid로 전환 \\
\texttt{wrong\_strategy}        & 모든 점수 $\approx 0$; 오라클 $\neq$ 현재 & 오라클 전략으로 전환 \\
\texttt{too\_narrow}            & 모든 점수 $\approx 0$; 이미 Hybrid & 쿼리 확장 \\
\texttt{missing\_entity}        & 쿼리 텍스트에 채워지지 않은 슬롯 & 슬롯 재채우기 \\
\texttt{insufficient\_coverage} & $0 < |\mathcal{R}| < \delta$; 현재 $\neq$ Hybrid & Hybrid로 전환 \\
\texttt{low\_quality}           & 일반적인 낮은 품질; 현재 $\neq$ Hybrid & Hybrid로 전환 \\
\bottomrule
\end{tabular}
\end{table}

\subsubsection*{답변 에이전트(AA).}

AA는 검증된 증거 청크 세트에서 텍스트 응답을 합성한다. 상위 $N$ 청크(기본 $N{=}10$, 검색 순위로 정렬)는 범주 레이블과 함께 번호가 매겨진 목록으로 형식화된다. LLM이 이용 불가능할 때는 상위 3개 증거 청크를 연결하는 추출 폴백을 사용한다.

\subsection{실험 구성}

4개의 실험이 3가지 주요 V2 설계 결정을 제거한다(표 \ref{tab:experiments}):

\begin{table}[htbp]
\centering
\caption{V2 제거 실험 구성}
\label{tab:experiments}
\begin{tabular}{lllll}
\toprule
\textbf{실험} & \textbf{명칭} & \textbf{QA 모드} & \textbf{점수화기} & \textbf{최대 루프} \\
\midrule
E1 & \texttt{v2\_oracle}            & 오라클    & 적응형    & 3 \\
E2 & \texttt{v2\_oracle\_noloop}    & 오라클    & 적응형    & 0 \\
E3 & \texttt{v2\_oracle\_rankproxy} & 오라클    & 순위 프록시 & 3 \\
E4 & \texttt{v2\_heuristic}         & 휴리스틱  & 적응형    & 3 \\
\bottomrule
\end{tabular}
\end{table}

\textbf{E1}(주요 V2 구성)은 오라클 추론 유형, 적응형 검증 점수화, 최대 3번의 재검색 루프를 사용하며 RQ8(V2 대 최고 V1)에 답한다. \textbf{E2}는 재검색을 비활성화하고(\texttt{MAX\_LOOPS}$=0$), E1--E2 차이가 재검색 루프에만 귀속되어 RQ9에 답한다. \textbf{E3}는 적응형 점수화기를 순위 프록시 점수화기로 대체하여 RQ10(적응형 점수화기 대 순위 프록시)에 답한다. \textbf{E4}는 오라클 분류를 휴리스틱 정규식 분류기로 교체하여 RQ11에 답한다.

% ─────────────────────────────────────────────────────────────────────
\section{결과}
\label{sec:results}

\subsection{V1 전체 성능}

표 \ref{tab:main}는 9개 시스템 모두에 대한 검색 지표를 보고한다. V1 실험 프로그램 전반에 걸쳐 몇 가지 핵심 정량적 관계가 나타난다. 하이브리드 RRF가 최고 고정 전략 성능을 달성한다: Recall@10$={0.745}$, MRR$={0.530}$, nDCG@10$={0.566}$. BM25에서 하이브리드로의 개선은 MRR +15.8\%로 연구에서 어느 시스템 전환보다 큰 절대 이득이다. 플래너 중에서, 규칙 플래너는 최고 Recall@10(0.755)과 Hit@10(0.808)을 달성하고, GPT-5는 최고 MRR(0.533)과 nDCG@10(0.570)을 달성한다.

\begin{table}[htbp]
\centering
\caption{전체 시스템 성능 비교 ($n{=}380$). 볼드체는 V1 시스템 중 열별 최고.}
\label{tab:main}
\begin{tabular}{lccccc}
\toprule
\textbf{시스템} & \textbf{Recall@10} & \textbf{MRR} & \textbf{nDCG@10} & \textbf{Hit@10} & \textbf{Recall@1} \\
\midrule
\multicolumn{6}{l}{\textit{V1 고정 전략 기준선}} \\
BM25               & 0.713 & 0.457 & 0.509 & 0.768 & 0.298 \\
밀집 (BGE-small)   & 0.740 & 0.501 & 0.548 & 0.787 & 0.346 \\
하이브리드 (RRF)   & 0.745 & 0.530 & 0.566 & 0.797 & 0.368 \\
\multicolumn{6}{l}{\textit{V1 적응형 플래너}} \\
규칙 플래너        & \textbf{0.755} & 0.522 & 0.563 & \textbf{0.808} & 0.358 \\
GPT-5 플래너       & 0.750 & \textbf{0.533} & \textbf{0.570} & 0.803 & \textbf{0.368} \\
\multicolumn{6}{l}{\textit{V2 에이전트 기반 RAG 실험}} \\
E1: 오라클 + 적응형 + 루프    & 0.743 & 0.529 & 0.566 & 0.797 & 0.368 \\
E2: 오라클 + 루프 없음        & 0.743 & 0.529 & 0.566 & 0.797 & 0.368 \\
E3: 오라클 + 순위 프록시 + 루프 & 0.743 & 0.529 & 0.566 & 0.797 & 0.368 \\
E4: 휴리스틱 + 적응형 + 루프  & 0.745 & 0.530 & 0.566 & 0.797 & 0.368 \\
\midrule
$\Delta$: E1 $-$ 규칙 플래너 (pp) & $-$1.2 & $+$0.7 & $+$0.3 & $-$1.1 & $+$1.0 \\
\bottomrule
\end{tabular}
\end{table}

\FloatBarrier

\subsection{추론 유형별 성능}

표 \ref{tab:by_type}는 V1 연구의 가장 중요한 발견을 드러낸다: 최고(집계: Recall@10$={0.924}$)와 최저(시계열: $0.000$--$0.200$) 추론 유형 사이의 50+ 퍼센트포인트 성능 격차. 이 격차는 쿼리 추론 유형이 검색 난이도의 지배적 예측 인자임을 확인한다.

\begin{table}[htbp]
\centering
\caption{추론 유형별 Recall@10: V1 시스템 및 V2 E1.}
\label{tab:by_type}
\begin{tabular}{lrrrrrr|c}
\toprule
\textbf{유형} & \textbf{n} & \textbf{BM25} & \textbf{밀집} & \textbf{하이브리드} & \textbf{규칙} & \textbf{GPT-5} & \textbf{E1 (V2)} \\
\midrule
집계     & 35  & 0.924 & 0.895 & 0.914 & 0.914 & 0.914 & 0.914 \\
비교     & 33  & 0.818 & 0.848 & 0.848 & 0.879 & 0.848 & 0.848 \\
단일\_도약 & 220 & 0.764 & 0.832 & 0.823 & 0.836 & 0.836 & 0.823 \\
다중\_도약 & 53  & 0.538 & 0.491 & 0.519 & 0.531 & 0.519 & 0.519 \\
예외     & 34  & 0.412 & 0.382 & 0.412 & 0.368 & 0.382 & 0.382 \\
시계열   & 5   & 0.200 & 0.000 & 0.100 & 0.200 & 0.100 & 0.200 \\
\midrule
\textbf{전체} & 380 & 0.713 & 0.740 & 0.745 & 0.755 & 0.750 & 0.743 \\
\bottomrule
\end{tabular}
\end{table}

\FloatBarrier

밀집 검색은 시계열 쿼리에서 유일하게 제로 재현율을 달성한다: 5개 시계열 쿼리가 집합적으로 제로 밀집-검색된 관련 청크를 받아, 날짜 토큰 순서가 정적 양방향 인코더 표현에 완전히 보이지 않음을 확인한다. 예외 쿼리는 Recall@10~$\leq$~0.412에서 보편적으로 정체하며, 어휘 및 의미적 신호 모두 조건부 절 의미를 캡처하지 못함을 확인한다.

\subsection{적응형 플래너 결과}

규칙 플래너는 6가지 추론 유형 중 4가지(시계열, 예외, 다중\_도약, 집계)에서 100\% SSA를 달성하여 최고 전체 Recall@10과 Hit@10을 달성한다. 표 \ref{tab:ssa}는 두 플래너에 대한 추론 유형별 전략 선택 정확도(SSA)를 분해한다.

\begin{table}[htbp]
\centering
\caption{추론 유형별 전략 선택 정확도(SSA)}
\label{tab:ssa}
\begin{tabular}{lrrcc}
\toprule
\textbf{추론 유형} & \textbf{n} & \textbf{오라클 수} & \textbf{규칙 SSA} & \textbf{GPT-5 SSA} \\
\midrule
시계열    & 5   & BM25(5) / Dense(0) / Hybrid(0) & 100.0\% & 0.0\%  \\
예외      & 34  & Dense(34)                       & 100.0\% & 82.4\% \\
다중\_도약 & 53  & Hybrid(53)                      & 100.0\% & 79.2\% \\
집계      & 35  & Hybrid(35)                      & 100.0\% & 77.1\% \\
비교      & 33  & Hybrid(33)                      & 100.0\% & 78.8\% \\
단일\_도약 & 220 & BM25(57) / Dense(34) / Hybrid(129) & 64.5\% & 79.1\% \\
\midrule
\textbf{전체} & 380 & BM25(62) / Dense(69) / Hybrid(249) & \textbf{78.4\%} & \textbf{77.6\%} \\
\bottomrule
\end{tabular}
\end{table}

GPT-5는 최고 전체 MRR(0.533)과 nDCG@10(0.570)을 달성하지만 시계열 쿼리에서 완전히 실패한다(0\% SSA: 5개 시계열 쿼리 모두 명시적 ``시계열$\to$bm25'' 지침에도 불구하고 Hybrid로 예측). GPT-5의 84개 전략 선택 오류의 질적 분석은 두 가지 지배적인 오류 모드를 드러낸다: API 문서 쿼리를 잘못 라우팅하는 \emph{범주 혼동}(55.3\%), GPT-5가 오라클 레이블과 일치하지 않는 방어 가능한 전략을 선택하는 \emph{오라클 불일치}(44.7\%).

\subsection{V2 제거 결과}

표 \ref{tab:ablation}는 쌍별 제거 결과를 보고한다. 각 실험 쌍은 하나의 설계 구성 요소를 분리하며, 다른 모든 것은 일정하게 유지된다.

\begin{table}[htbp]
\centering
\caption{V2 제거: 증분적 구성 요소 기여.}
\label{tab:ablation}
\begin{tabular}{lllccc}
\toprule
\textbf{쌍} & \textbf{제거된 구성 요소} & \textbf{RQ} & \textbf{$\Delta$ R@10} & \textbf{$\Delta$ MRR} & \textbf{$\Delta$ nDCG@10} \\
\midrule
E1 $-$ E2 & 재검색 루프 (3 루프 대 0) & RQ9  & $0.0$ & $0.0$ & $0.0$ \\
E1 $-$ E3 & 적응형 점수화기 대 순위 프록시 & RQ10 & $0.0$ & $0.0$ & $0.0$ \\
E1 $-$ E4 & 오라클 QA 대 휴리스틱 QA   & RQ11 & $-0.1$ & $0.0$ & $-0.1$ \\
\bottomrule
\end{tabular}
\end{table}

표 \ref{tab:pipeline}는 V2 시스템의 내부 운영을 반영하는 파이프라인 통계를 보고한다.

\begin{table}[htbp]
\centering
\caption{V2 파이프라인 통계: 검증 통과율 및 루프 분포 ($n{=}380$).}
\label{tab:pipeline}
\begin{tabular}{lcccc}
\toprule
\textbf{실험} & \textbf{통과율} & \textbf{평균 루프} & \textbf{0-루프 (\%)} & \textbf{$\geq$3 루프 (\%)} \\
\midrule
E1: 오라클 + 적응형 + 루프    & 98.7\% & 0.039 & 98.7\% & 1.3\% \\
E2: 오라클 + 루프 없음        & 98.7\% & 0.000 & 100\%  & ---  \\
E3: 오라클 + 순위 프록시 + 루프 & 100\% & 0.000 & 100\%  & 0\%  \\
E4: 휴리스틱 + 적응형 + 루프  & 97.6\% & 0.071 & 97.6\% & 2.4\% \\
\bottomrule
\end{tabular}
\end{table}

% ─────────────────────────────────────────────────────────────────────
\section{토론}
\label{sec:discussion}

\subsection{시스템 간 비용-성능 상충관계}

5가지 V1 시스템은 운영 비용과 지연에서 5차원 범위에 걸쳐 있다. BM25와 밀집 검색은 인덱스 구성 후 제로 한계 비용으로 밀리초 미만의 지연으로 운영된다. 하이브리드 RRF는 무시 가능한 오버헤드를 추가한다. 규칙 플래너는 쿼리당 0.1~ms 미만을 기여한다. GPT-5는 평균 지연 5,649~ms와 \$0.016/쿼리의 LLM API 호출을 도입한다.

비용-성능 경계에서, 규칙 플래너가 지배한다: 효과적으로 제로 운영 비용과 밀리초 미만의 지연으로 최고 Recall@10(0.755)을 달성한다. GPT-5는 세 자릿수 더 높은 지연과 정량화 가능한 금전적 비용으로 순위 품질 개선(MRR +2.0\%, nDCG@10 +1.2\%)을 제공한다. V2 파이프라인의 검색 비용 증분은 무시 가능하다: 쿼리당 평균 재검색 루프 0.039(E1)와 0.071(E4)(표 \ref{tab:pipeline})로, 평균 쿼리는 0.1번의 추가 검색 호출 미만을 트리거한다. V2의 총 추가 검색 비용은 기준 검색 비용의 약 4--7\%로, 측정 노이즈 수준 이하이다.

\subsection{검색 아키텍처가 주요 성능 동인}

전체 실험 프로그램에 걸쳐 가장 일관된 발견은 검색 아키텍처 선택이 계획 메커니즘 선택보다 성능을 훨씬 더 많이 이끈다는 것이다. BM25에서 하이브리드 RRF로의 개선(MRR +15.8\%)은 하이브리드 기준선 대비 최고 플래너 개선(GPT-5의 MRR +2.0\%)보다 8배 크다. 이 비율은 중요한 실용적 함의를 갖는다: 기업 RAG 설계자는 검색 아키텍처 선택을 일차 설계 결정으로 우선시하고, 검색 인프라가 올바르게 구성된 후에만 계획 지능에 투자해야 한다.

이 발견은 특히 RAG 문헌의 일반적인 직관과 상반되기 때문에 중요하다: 지능적인 쿼리 라우팅 또는 LLM 기반 쿼리 이해가 검색 품질 향상의 주요 레버라는 직관. 우리의 결과는 반대 순서를 정량화한다. 설명은 기업 코퍼스의 특성에 있다: SEKD가 보완적인 신호 요구사항을 가진 6가지 문서 범주를 포함하기 때문에 BM25와 밀집 신호 모두 포괄적인 커버리지를 위해 필요하다.

\subsection{GPT-5 대비 규칙 기반 계획이 경쟁력 있는 이유}

규칙 플래너(SSA$={78.4\%}$, Recall@10$={0.755}$)와 GPT-5(SSA$={77.6\%}$, Recall@10$={0.750}$)의 경험적 근사 동등성은 명시적 메타데이터를 가진 구조화된 계획 작업에서 프론티어 LLM이 실질적인 이점을 제공한다는 가정에 도전한다.

세 가지 요인이 규칙 플래너의 경쟁력 있는 성능을 설명한다. \textit{오라클-규칙 구조적 정렬}: 오라클 레이블이 규칙 플래너의 규칙에 인코딩된 동일한 추론 유형-전략 매핑을 사용하여 할당되었기 때문에, 평가는 규칙 플래너에 구조적으로 유리하다. \textit{낮은 작업 복잡도}: 명시적 구조화된 메타데이터에 기반한 세 가지 범주에 걸친 검색 전략 분류는 LLM 추론 능력을 활용하지 않는 낮은 복잡도 작업이다. \textit{보편적으로 올바른 기본값으로서의 Hybrid}: 65.5\%의 오라클 레이블이 Hybrid이므로, 불확실할 때 Hybrid를 기본값으로 사용하는 모든 시스템은 특정 지식 없이도 최소 65.5\% SSA를 달성한다.

GPT-5의 시계열 쿼리에 대한 완전한 실패(명시적 ``시계열$\to$bm25'' 프롬프트 지침에도 불구하고 0\% SSA)는 연구에서 가장 이론적으로 중요한 부정적 결과이다. 이는 프론티어 LLM이 전 세계적으로 일반적인 패턴과 충돌할 때 명시적으로 명시된 계획 규칙을 체계적으로 무시할 수 있음을 입증한다. 안전이 중요한 라우팅 결정에서, 결정론적 규칙 집행이 LLM 명령 따르기보다 더 신뢰할 수 있다.

\subsection{예외 및 시계열 쿼리: 구조적 실패 모드}

V1에서 쉬운(집계: 0.924)과 어려운(시계열: $\leq 0.200$) 추론 유형 사이의 50+ 퍼센트포인트 성능 격차는 평가된 시스템의 매개변수적 결함에 귀속될 수 없다. 실패는 구조적이다: 쿼리 유형과 이용 가능한 검색 신호 공간 사이의 불일치를 반영한다.

\textbf{시계열 쿼리}는 정책 버전, API 릴리스, 회의 결정 간의 날짜 순서 순차적 관계를 대상으로 한다. 밀집 검색은 날짜 토큰 순서가 정적 양방향 인코더 표현에 완전히 보이지 않아 시계열 쿼리에서 유일하게 제로 재현율을 달성한다. BM25는 정확한 날짜 문자열 매칭이 때때로 관련 시계열 참조를 포함하는 특정 청크를 검색하기 때문에 0.200을 달성하지만, 이는 체계적이 아닌 우연적이다. 시계열 검색 실패에 대한 솔루션은 검색 전략 선택을 넘어서는 아키텍처적 변경이 필요하다: BM25에서 날짜 필드 부스팅, 밀집 모델에서 시계열 위치 인코더, 또는 명시적 시계열 메타데이터 인덱싱.

\textbf{예외 절 쿼리}는 정책 텍스트의 조건부 수정자를 대상으로 한다(``직원이 면제되는 조건은?''). BM25는 쿼리의 예외 조건에 설명된 예외 절 언어가 키워드 중복되지 않기 때문에 실패한다. 밀집 검색은 일반 텍스트로 학습된 의미적 인코더가 청크 수준에서 예외 절을 일반 정책 문장과 구별하지 못하기 때문에 실패한다. 올바른 솔루션은 정책 예외 절 쌍에 대한 대조적 미세 조정이지만, 이는 본 연구 범위 밖이다.

이러한 실패 모드가 V2 에이전트 확장의 동기가 되었으며, 반복적 검증 및 전략 전환을 통해 이를 해결하도록 설계되었다. 그러나 V2 평가가 드러내듯이, 검증 임계값이 너무 허용적임이 밝혀졌다: 쿼리의 98.7\%가 첫 번째 검색 시도에서 통과하여(표 \ref{tab:pipeline}), 재검색 루프가 380개 쿼리 중 단 5개에 대해서만 발동된다. 따라서 V2는 이러한 구조적 실패 모드에서 측정 가능한 개선을 생성하지 않는다.

\subsection{V2 에이전트 파이프라인: 제거 결과 해석}

제거 결과(표 \ref{tab:ablation})는 3가지 구성 요소 대조 전반에 걸쳐 균일한 귀무 결과를 드러낸다. E1--E2 비교(재검색 루프)는 $\Delta$(Recall@10)$=0.0$, $\Delta$(MRR)$=0.0$, $\Delta$(nDCG@10)$=0.0$을 산출한다. 루프는 380개 쿼리 중 단 5개(1.3\%)에 대해서만 발동되며, 모두 \texttt{insufficient\_coverage} 실패로 분류된다. 적응형 점수화기가 98.7\%의 쿼리를 첫 번째 시도에서 통과시키기 때문에, 루프의 수정 기여는 소수점 두 자리에서 반올림 정밀도 이하이다.

E1--E3 비교(적응형 점수화기 대 순위 프록시)도 세 가지 지표 전반에 걸쳐 0.0을 산출한다. E3의 순위 프록시 점수화기는 100\% 첫 번째 패스 검증 통과율을 달성하여(표 \ref{tab:pipeline}) 재검색 루프를 완전히 불가능하게 만든다. E1--E4 비교(오라클 대 휴리스틱 QA)는 유일한 비제로 델타를 생성한다: E1 대 E4에 대해 $-0.1$~pp Recall@10 및 $-0.1$ nDCG@10. 이 반직관적 결과---오라클 QA가 휴리스틱 QA를 \textit{하회하는}---는 구조적 비대칭을 반영한다. 휴리스틱 QAA는 어떤 규칙도 매칭되지 않을 때 Hybrid 검색을 기본값으로 사용하며, Hybrid(Recall@10$={0.745}$)가 두 번째로 최고 V1 시스템이다. 오라클 규칙 플래너는 예외 쿼리를 Dense(Recall@10$={0.382}$)로 라우팅하는데, 기술적으로 오라클 레이블에 의해 올바르지만 실제로는 그 쿼리에서 Hybrid보다 낮은 성능을 보인다.

지배적인 진단 발견은 임계값 보정 실패이다. 검증 통과 조건(수식 \ref{eq:pass})은 주제에 인접하지만 쿼리가 필요로 하는 특정 정책 절이나 시계열 참조를 포함하지 않는 청크가 임계값 이상의 점수를 받는 경우에도 임계값 이상으로 계산될 만큼 너무 허용적인 고정 유형별 임계값 $\theta_\tau$를 사용한다. 비 LLM 점수화기는 표현 수준에서 ``근접하지만 잘못된'' 청크와 ``올바른'' 청크를 구별할 수 없어, 수정 루프 반복을 트리거해야 하는 예외 및 시계열 쿼리에 대한 허위 검증 통과를 생성한다.

\subsection{기업 RAG 시스템 설계를 위한 함의}

우리의 발견은 기업 RAG 실무자를 위한 6가지 구체적이고 정량적으로 근거한 설계 권고사항을 도출한다.

\textbf{(D1) 하이브리드 RRF로 시작하라.} 하이브리드 RRF 기준선은 학습이 필요 없고, BM25에 비해 MRR +15.8\% 개선을 제공하며, 최고 계획 개선보다 8배 큰 수익을 제공한다. 이것이 검색 스택에서 단일 최고 수익 투자이다.

\textbf{(D2) 무료 개선으로 규칙 플래너를 사용하라.} 7개 결정론적 규칙이 제로 운영 비용으로 Hybrid 대비 Recall@10 +1.0 pp를 추가한다. 플래너는 감사 가능하고, 해석 가능하며, 지연 중립적이다.

\textbf{(D3) 정밀도 중요 저용량 애플리케이션에만 GPT-5 계획을 배포하라.} 상위 1 또는 상위 3 결과 품질이 다운스트림 결정 품질을 결정하는 애플리케이션에서는 MRR +2.0\% 이점이 의미 있다. 고용량 검색, 배치 처리, 또는 \$0.016/쿼리가 중요한 상황에서는 규칙 플래너가 강력하게 선호된다.

\textbf{(D4) 계획이 아닌 인프라를 통해 시계열 및 예외 쿼리를 해결하라.} 어떤 계획이나 에이전트 메커니즘도 쿼리 유형과 검색 신호 사이의 근본적인 불일치를 보완할 수 없다. 시계열 쿼리는 구조화된 BM25 인덱스에서 날짜 필드 부스팅이나 전용 시계열 지식 그래프 레이어가 필요하다. 예외 절 쿼리는 정책 예외 쌍에 대한 대조적 미세 조정이나 인덱스 시간에 예외 절 영역을 식별하는 분류기가 필요하다.

\textbf{(D5) 다중 도약 쿼리에 에이전트 기반 재검색을 사용하라.} 다중 도약 쿼리(모든 V1 시스템에서 Recall@10~$\leq$~0.538)는 단일 패스 검색이 충족할 수 없는 순차적 증거 종속성을 포함한다. V2 파이프라인의 슬롯 채우기를 포함하는 순차적 검색 모드가 이 구조적 제한을 해결한다.

\textbf{(D6) 검색 분포 이동을 감지하기 위해 검증 통과율을 모니터링하라.} V2 검증 에이전트의 통과율은 검색 품질 모니터링을 위한 자연스러운 운영 지표이다. 생산 쿼리 스트림에서 감소하는 통과율은 코퍼스 드리프트, 임베딩 분포 이동, 또는 플래너의 규칙으로 다루지 않는 새로운 쿼리 유형의 출현을 나타내며, 검색 구성 재평가를 트리거해야 한다.

\subsection{한계 및 타당성 위협}

\textbf{합성 코퍼스.} SEKD 문서는 실제 기업 코퍼스의 언어적 다양성, 문서 간 상호 참조, 도메인별 전문 용어가 부족하다. SEKD 발견의 실제 기업 배포에 대한 일반화는 생산 문서 세트에 대한 검증이 필요하다.

\textbf{오라클 레이블 품질.} 휴리스틱으로 할당된 오라클 레이블은 규칙 플래너를 구조적으로 선호한다. 경험적으로 도출된 오라클 레이블(쿼리당 Recall@10을 최대화하는 전략)이 더 객관적인 평가 기준선을 제공하고, GPT-5 오류의 44.7\%가 \texttt{oracle\_disagreement}로 분류된 모호한 불일치를 해결할 것이다.

\textbf{쿼리 볼륨 및 통계적 유의성.} 전체 쿼리 세트(380개 답변 가능 쿼리)는 95\% 신뢰 수준에서 $\geq$1.5 pp의 집계 지표 차이에 적절한 통계적 검정력을 제공하지만, 하위 그룹 분석(특히 $n{=}5$인 시계열 추론)은 점 추정 정밀도에 충분한 검정력이 부족하다. 5개의 시계열 쿼리는 통계적 유의성에 불충분하다. 시계열 발견은 정확한 점 추정이 아닌 방향성으로 해석되어야 한다.

\textbf{단일 임베딩 모델.} BGE-small-en-v1.5만 평가되었다. 더 큰 밀집 모델(text-embedding-3-large, E5-large-v2)은 특히 의미적 인코딩 품질이 제한 요인인 예외 및 시계열 쿼리에서 상당히 다른 순위를 생성할 수 있다.

\textbf{V2 답변 품질.} 본 논문은 검색 지표에서만 V2를 평가한다. 답변 품질(사실적 정확도, 인용 정밀도, 응답 완전성)은 평가되지 않았으며, 향후 연구를 위한 중요한 미해결 문제를 구성한다.

% ─────────────────────────────────────────────────────────────────────
\section{결론}
\label{sec:conclusion}

우리는 목적에 맞게 구축된 평가 코퍼스(SEKD), 5가지 시스템 V1 기준선 연구, 4개의 통제된 제거 실험을 포함하는 5-에이전트 V2 에이전트 파이프라인을 결합하여 기업 지식 관리를 위한 적응형 검색 계획과 에이전트 기반 재검색에 대한 포괄적인 실험 연구를 제시하였다. V1 연구는 세 가지 주요 발견을 생성하였다. 첫째, 하이브리드 역순위 융합이 이질적인 기업 코퍼스를 위한 올바른 검색 기반이다: BM25 단독 대비 MRR +15.8\% 개선이 최고 계획 개선보다 8배 크며, 검색 아키텍처가 주요 성능 동인임을 확인한다. 둘째, 7개 규칙 결정론적 플래너가 제로 운영 비용과 밀리초 미만 지연으로 전체 V1 시스템 중 최고 Recall@10($0.755$)을 달성하며, GPT-5가 \$0.016/쿼리의 비용으로 한계적 순위 품질 이점(MRR +2.0\%)을 제공한다. 셋째, 예외 절 쿼리(Recall@10~$\leq$~0.412)와 시계열 쿼리(Recall@10~$\leq$~0.200)는 쿼리 유형과 검색 신호 공간 사이의 근본적인 불일치로 인해 5가지 V1 시스템 모두 하에서 불가능한 구조적 실패 모드를 나타낸다.

이러한 실패에 동기를 받아, V2 에이전트 RAG 파이프라인은 반복적 검증과 재검색을 도입한다. 4개의 제거 실험(E1--E4)은 재검색 루프(E1 대 E2), 적응형 검증 점수화 대 순위 프록시 점수화(E1 대 E3), 오라클 대 휴리스틱 쿼리 분류(E1 대 E4)의 기여를 분리하도록 설계된다. 결과는 균일한 귀무 결과를 드러낸다: 재검색 루프는 유형별 임계값이 너무 허용적이어서 예외 및 시계열 쿼리가 첫 번째 검색에서 검증을 통과하여 재검색 루프를 억제하기 때문에 전혀 발동되지 않는다.

기업 RAG 실무자를 위한 완전한 실험 프로그램은 명확한 우선순위 스택을 생성한다: (1)~하이브리드 RRF를 검색 기반으로 수립; (2)~대상 전략 라우팅을 위한 규칙 기반 플래너 레이어 추가; (3)~상위 1 또는 상위 3 결과에 대한 순위 품질이 중요하고 쿼리 볼륨이 낮을 때만 GPT-5 계획 배포; (4)~계획이 아닌 검색 인프라 변경을 통해 예외 및 시계열 실패 해결; (5)~다중 도약 쿼리 유형이 워크로드의 상당 부분을 차지할 때 V2 에이전트 파이프라인 배포.

\paragraph{향후 연구.} 식별된 한계에서 여러 연구 방향이 도출된다. 첫째, 경험적으로 도출된 오라클 레이블(열거로 쿼리당 Recall@10을 최대화하는 전략)이 더 객관적인 플래너 평가 기준선을 제공하고 오라클 불일치 모호성을 해결할 것이다. 둘째, 날짜 필드 부스팅, 구조화된 시계열 메타데이터 추출, 또는 전용 타임라인 레이어를 포함하는 시계열 인식 인덱싱이 상당한 시계열 추론 요구사항을 가진 기업 코퍼스를 위한 최우선 인프라 투자이다. 셋째, 정책 예외 절 쌍에 대한 밀집 인코더의 대조적 미세 조정이 예외 절 검색 실패 모드를 직접 해결할 것이다. 넷째, 쿼리별 루브릭을 가진 LLM 판사 프로토콜을 사용한 V2 답변 에이전트의 응답 품질(사실적 정확도, 인용 정밀도, 완전성) 평가가 필요하다. 다섯째, 사용자 피드백 신호(클릭스루, 쿼리 재구성, 세션 길이)로부터의 온라인 학습이 완전한 재평가 없이 배포에서 규칙 플래너의 라우팅 결정을 적응시킬 수 있다. 마지막으로, SEKD를 합성 문서와 함께 실제 기업 문서를 포함하도록 확장하면 여기서 보고된 발견에 대한 중요한 외부 타당성 확인을 제공할 것이다.

% ─────────────────────────────────────────────────────────────────────
\bibliography{references}

\end{document}
